\section{可微几何优化与端到端 SfM}

\subsection{从模块增强到系统级重构}

如果说对应学习解决的是前端观测质量问题，那么可微 SfM 则进一步触及传统多视图几何的核心系统结构。过去许多深度学习方法只替换匹配、深度或位姿初始化等局部模块，而后端仍沿用不可微的经典优化器。这使得系统难以真正实现端到端学习，也限制了最终误差对前端表征的反馈能力\cite{wang2021deep2view}。

\textit{VGGSfM: Visual Geometry Grounded Deep Structure From Motion} 在这方面迈出关键一步\cite{wang2024vggsfm}。该工作将深度点轨迹、相机恢复、三角化和 Bundle Adjustment 组织到同一可微框架中，使传统 SfM 系统中的多个关键环节能够在统一目标下训练。与仅对某个部件做学习增强的方案相比，VGGSfM 更接近对整个结构恢复问题的“系统级重构”。

\subsection{VGGSfM 的核心思想}

VGGSfM 的一个重要特点是同时强调“几何 grounded”与“深度学习驱动”。其输入不再是传统关键点匹配链，而是由点轨迹提供的更稳定观测；其相机恢复不是增量式逐步注册，而是尝试同步估计全部相机；其三维点优化则通过可微的 BA 层来完成\cite{wang2024vggsfm}。这意味着模型不仅在前端表征上使用学习方法，也在后端优化层面尽量保留传统几何问题的可解释形式。

从研究主线角度看，VGGSfM 是 VGG 团队从“增强几何管线”走向“重构几何管线”的转折点。一方面，它延续了前一阶段对轨迹、匹配和几何约束的重视；另一方面，它通过可微化方式将原本分离的多个模块组织到统一训练框架中，为进一步走向前馈式统一几何预测创造了条件。

\subsection{与传统 SfM 及同类方法的关系}

相较于 COLMAP 等传统系统，VGGSfM 的优势主要在于能够借助学习表征改善鲁棒性，并通过统一训练强化全局一致性\cite{schoenberger2016colmap,wang2024vggsfm}。相较于只在前端使用深度特征、后端仍完全依赖传统不可微求解器的混合方案，VGGSfM 则体现出更强的系统整合能力。

不过，VGGSfM 仍然保留了优化层，说明这一阶段的重点仍是“让几何优化可学习”，而非完全放弃优化。这一点在研究路径上十分关键：它既保留了传统几何的确定性和解释性，也暴露出后续研究进一步追求效率与统一性的方向，即能否在不依赖复杂测试时优化的前提下，直接输出关键三维属性。
